\documentclass[11pt]{article}

\usepackage[utf8]{inputenc}
\usepackage[T1]{fontenc}
\usepackage{geometry}
\usepackage{amsmath}
\usepackage{amssymb}
\usepackage{booktabs}
\usepackage{hyperref}
\usepackage{xurl}
\usepackage{xcolor}
\usepackage{listings}
\usepackage{enumitem}
\usepackage{parskip}
\usepackage{graphicx}
\graphicspath{{../images/}}

\geometry{margin=1in}

\lstset{
  language=Java,
  basicstyle=\ttfamily\small,
  keywordstyle=\color{blue},
  commentstyle=\color{gray},
  stringstyle=\color{teal},
  showstringspaces=false,
  columns=fullflexible,
  frame=single,
  framerule=0.2pt,
  breaklines=true
}

\setlist[itemize]{topsep=0.3em,itemsep=0.2em}
\setlist[enumerate]{topsep=0.3em,itemsep=0.2em}

% Simple per-section source line.
\newcommand{\notesources}[1]{\par\noindent{\small\color{gray}\textit{Sources:} \sloppy #1\par}}

% Unit-wise source strings for this subject (used by slides/notes).
% Path is relative to the lecture `latex/` folder (the compilation CWD).
% OOPS with Java (CSEG1044) — unit-wise sources
%
% These are short, student-facing reference strings intended to be used with:
%   - \slidesources{...}  (in beamer slides)
%   - \notesources{...}   (in lecture notes)
%
% Style inspiration: other subjects in My-Subjects/ use semicolon-separated sources
% and include an "accessed" date for web documentation.

\newcommand{\OOPSUnitISources}{Schildt, \textit{Java: A Beginner's Guide} (10e), McGraw-Hill (Java fundamentals + OOP basics).; Oracle Java Tutorials: Getting Started + Language Basics + Classes and Objects (accessed \today).; Java SE API docs (e.g., \texttt{String}, \texttt{Scanner}) (accessed \today).}

\newcommand{\OOPSUnitIISources}{Schildt, \textit{Java: The Complete Reference} (12e), McGraw-Hill (classes, inheritance, packages, interfaces).; Bloch, \textit{Effective Java} (3e), Addison-Wesley, 2018 (design choices; interfaces vs abstract classes).; Oracle Java Tutorials: Classes and Objects + Interfaces + Packages (accessed \today).; Java SE API docs (e.g., \texttt{Comparable}, \texttt{Runnable}) (accessed \today).}

\newcommand{\OOPSUnitIIISources}{Schildt, \textit{Java: The Complete Reference} (12e) (nested/inner classes, exceptions, threads, I/O).; Oracle Java Tutorials: Nested Classes + Exceptions + Concurrency + Basic I/O (accessed \today).; Java SE API docs (e.g., \texttt{Thread}, \texttt{AutoCloseable}) (accessed \today).}

\newcommand{\OOPSUnitIVSources}{Schildt, \textit{Java: The Complete Reference} (12e) (generics, lambdas, Swing, JDBC).; Bloch, \textit{Effective Java} (3e) (generics + lambdas).; Oracle Java Tutorials: Generics + Lambda Expressions + Swing + JDBC Basics (accessed \today).; Java SE API docs (e.g., \texttt{java.util.function}, \texttt{javax.swing}, \texttt{java.sql}) (accessed \today).}

\newcommand{\OOPSUnitVSources}{Schildt, \textit{Java: The Complete Reference} (12e) (Collections Framework + wrapper classes).; Oracle Java docs: Collections Framework overview (accessed \today).; Java SE API docs (\texttt{java.util} collections; wrapper classes in \texttt{java.lang}) (accessed \today).}

\newcommand{\OOPSUnitVISources}{Git SCM: \textit{Pro Git} (2e) (version control workflow), accessed \today.; GitHub Docs (repositories, issues, pull requests), accessed \today.; Oracle Java Tutorials: Swing + JDBC (accessed \today).; JUnit 5 User Guide (accessed \today).; Apache Maven Project: Getting Started (accessed \today).; Gradle User Manual: Getting Started (accessed \today).; UML class diagrams: UML 2.x overview/spec (accessed \today).}

\newcommand{\OOPSMiscWhyInterfacesSources}{Schildt, \textit{Java: The Complete Reference} (12e) (interfaces).; Bloch, \textit{Effective Java} (3e) (prefer interfaces where appropriate).; Oracle Java Tutorials: Interfaces (accessed \today).; Java SE API docs (e.g., \texttt{Comparable}, \texttt{Runnable}) (accessed \today).}



\title{Object Oriented Programming (CSEG1044)\\Unit 05 -- Lecture 03 Notes\\Wrapped Collections + `Collections` Utilities + Wrappers}
\begin{document}
\maketitle
\tableofcontents

\section{Lecture Snapshot}
\begin{itemize}
  \item \textbf{Unit focus:} Collections Framework and Wrapper Classes
  \item \textbf{Main new ideas:} \texttt{Collections} utilities (common ones), Wrapped collections (why/when), Wrapper classes + autoboxing pitfalls
  \item \textbf{Course thread:} Use Wrapped Collections + `Collections` Utilities + Wrappers to manage in-memory project data more cleanly and predictably inside Campus Resource Manager.
\end{itemize}
\notesources{\OOPSUnitVSources}

\section{Lecture Overview}
This lecture collects three practical topics that show up in day-to-day Java code:
\begin{itemize}
  \item utility methods in \texttt{java.util.Collections},
  \item wrapper ``views'' around collections (unmodifiable, synchronized, checked),
  \item wrapper classes (\texttt{Integer}, \texttt{Double}, ...) and autoboxing pitfalls.
\end{itemize}
We also briefly introduce the class loading idea from the syllabus.
\notesources{\OOPSUnitVSources}

\section{Prerequisite Bridge}
Before reading the new material in depth, do a fast recall of the older ideas listed below.
\begin{itemize}
  \item \textbf{Collection interfaces}: Utilities and wrappers only make sense once List, Set, and Map roles are already clear. Main reference: \texttt{Unit 05 / Lecture 01 slides/notes}.
  \item \textbf{Primitive and object types}: Wrapper classes connect directly to the type-conversion ideas from Unit 01. Main reference: \texttt{Unit 01 / Lecture 03 slides/notes}.
\end{itemize}
This lecture only revisits those ideas briefly so that the main explanation can stay focused on the new concept sequence.
\notesources{\OOPSUnitVSources}

\section{Learning Outcomes}
After this lecture, you should be able to:
\begin{itemize}
  \item Use common \texttt{Collections} utilities (sort, min/max, shuffle, reverse).
  \item Use wrapped collections (unmodifiable/synchronized/checked) appropriately.
  \item Avoid common wrapper/autoboxing pitfalls (and know class loading basics).
\end{itemize}

\section{Suggested Reading Order}
\begin{enumerate}
  \item Read the \textbf{Prerequisite Bridge} first and confirm each recalled idea.
  \item Study the central explanation in this order: \textit{\texttt{Collections} utilities (common ones)} $\rightarrow$ \textit{Wrapped collections (why/when)} $\rightarrow$ \textit{Wrapper classes + autoboxing pitfalls} $\rightarrow$ \textit{Class loading basics (intro)}.
  \item Trace the worked example line by line before checking short answers.
  \item Attempt the practice section without looking at the solution immediately.
  \item Finish with the exit questions and further-study suggestions to confirm retention.
\end{enumerate}
\notesources{\OOPSUnitVSources}

\section{Key Terms (Quick)}
\begin{itemize}
  \item \textbf{Utility class}: class with static helper methods (e.g., \texttt{Collections}).
  \item \textbf{Unmodifiable view}: wrapper that blocks structural modification.
  \item \textbf{Synchronized wrapper}: wrapper that adds basic synchronization around operations.
  \item \textbf{Checked wrapper}: wrapper that enforces runtime type checks on adds.
  \item \textbf{Wrapper class}: object form of a primitive (e.g., \texttt{Integer} for \texttt{int}).
  \item \textbf{Autoboxing/unboxing}: automatic conversion between primitive and wrapper.
  \item \textbf{Class loading}: JVM loads class bytecode into memory when needed.
\end{itemize}

\section{Detailed Notes}
\subsection{\texttt{Collections} utilities (common ones)}
\texttt{java.util.Collections} provides algorithms for collections:
\begin{itemize}
  \item \texttt{sort(list)} or \texttt{sort(list, comparator)}
  \item \texttt{reverse(list)}
  \item \texttt{shuffle(list)}
  \item \texttt{min(list)}, \texttt{max(list)}
  \item \texttt{frequency(collection, element)}
  \item \texttt{binarySearch(list, key)} (requires sorted list)
\end{itemize}

\subsection{Wrapped collections (why/when)}
Wrapped collections are ``views'' (wrappers) around an existing collection.
They are used to enforce constraints without copying:
\begin{itemize}
  \item \texttt{unmodifiableList}: when you want to expose read-only data to callers.
  \item \texttt{synchronizedList}: basic thread-safety wrapper (older style; modern code may use concurrent collections).
  \item \texttt{checkedList}: runtime type checking (useful when mixing legacy/raw types).
\end{itemize}

\textbf{Important:} ``unmodifiable'' means you cannot change structure via the wrapper, but if the underlying list changes, the view reflects it.

\subsection{Wrapper classes + autoboxing pitfalls}
Wrapper classes are needed because generics work with reference types, not primitives:
\begin{itemize}
  \item \texttt{List<Integer>} not \texttt{List<int>}
\end{itemize}

\textbf{Pitfalls:}
\begin{itemize}
  \item Unboxing null $\rightarrow$ \texttt{NullPointerException}\par
  Example: \texttt{Integer x = null; int y = x;} crashes.
  \item Using \texttt{==} for wrapper comparisons can be confusing.\par
  Prefer \texttt{equals()} or compare primitives.
  \item Performance: excessive boxing/unboxing in tight loops can slow code.
\end{itemize}

\subsection{Class loading basics (intro)}
The JVM loads classes when they are needed (first time referenced).
\path|Class.forName("com.mysql.cj.jdbc.Driver")| is sometimes used to load a JDBC driver explicitly (older pattern).
If a class is missing, you may see \texttt{ClassNotFoundException}.
\notesources{\OOPSUnitVSources}

\section{Running Project Connection}
Use Wrapped Collections + `Collections` Utilities + Wrappers to manage in-memory project data more cleanly and predictably inside Campus Resource Manager.
\begin{itemize}
  \item Use Campus Resource Manager as the running case study, or map the same idea to your own project domain if you are using another example.
  \item Ask where today's idea should live: model, service, DAO, UI, or team workflow.
  \item Use the lecture example as a smaller version of the same design choice you will make in the capstone.
\end{itemize}
\notesources{\OOPSUnitVSources}

\section{Common Mistakes and Confusions}
\begin{itemize}
  \item Assuming wrapper objects behave exactly like primitives in every context.
  \item Treating synchronized wrappers as a full concurrency design.
  \item Forgetting that an unmodifiable view can still reflect external changes to the backing collection.
\end{itemize}
\notesources{\OOPSUnitVSources}

\section{Worked Example: utilities + wrappers + autoboxing}
\begin{lstlisting}
import java.util.*;

public class Demo {
    public static void main(String[] args) {
        List<Integer> a = new ArrayList<>(Arrays.asList(3, 1, 2));
        Collections.sort(a);
        System.out.println("Sorted: " + a);
        System.out.println("Min=" + Collections.min(a) + ", Max=" + Collections.max(a));

        List<Integer> ro = Collections.unmodifiableList(a);
        // ro.add(99); // UnsupportedOperationException

        List checked = Collections.checkedList(new ArrayList(), String.class);
        checked.add("ok");
        // checked.add(123); // ClassCastException at add-time

        Integer x = null;
        // int y = x; // NullPointerException (unboxing null)
    }
}
\end{lstlisting}

\section{Demo / Observation Task}
Use this block as a short reproducible demo or observation task.
\begin{itemize}
  \item Reproduce the lecture's main example around \textit{\texttt{Collections} utilities (common ones)} once without looking at the solution first.
  \item Change one input, rule, or design choice in Campus Resource Manager and predict the result before checking it.
  \item Record one success case, one edge case, and one failure case so the idea becomes observable instead of purely theoretical.
\end{itemize}
\notesources{\OOPSUnitVSources}

\section{Practice Check (with brief solutions)}
\begin{enumerate}
  \item What exception occurs if you try to add to an unmodifiable list?\par
  \textbf{Answer:} \texttt{UnsupportedOperationException}.

  \item Why do we need wrapper classes with collections?\par
  \textbf{Answer:} generics require reference types; primitives cannot be used directly.

  \item What happens if you unbox \texttt{Integer x = null}?\par
  \textbf{Answer:} \texttt{NullPointerException}.
\end{enumerate}

\section{Self-Study Practice (no solutions)}
\begin{enumerate}
  \item Shuffle a list of integers and print the result.
  \item Create an unmodifiable view of a list and test which operations fail.
  \item Write a program that reads integers into a \texttt{List<Integer>} and then finds min/max using \texttt{Collections}.
\end{enumerate}

\section{Lab Alignment (recommended)}
\begin{itemize}
  \item Exp-18: Collections Framework
  \item Exp-19: Wrapper Classes
\end{itemize}

\section{Further Study}
\begin{itemize}
  \item Revisit the prerequisite references if any recall bullet still feels weak.
  \item Rework the lecture example with one changed assumption, input, or design choice.
  \item Read the textbook and documentation references listed in the source line for a second explanation of the same topic.
  \item Apply the same idea once inside Campus Resource Manager so that the concept moves from theory to design practice.
\end{itemize}
\notesources{\OOPSUnitVSources}

\section{Exit Questions (with sample answers)}
\textbf{Q1.} What is the difference between \texttt{Collections} and \texttt{Collection}?\par
\textbf{Sample answer:} \texttt{Collection} is an interface representing a group of elements. \texttt{Collections} is a utility class containing static helper methods for working with collections.

\vspace{0.6em}
\textbf{Q2.} Why can unboxing a null \texttt{Integer} crash?\par
\textbf{Sample answer:} If \texttt{x} is null, converting it to a primitive triggers \texttt{NullPointerException}.

\vspace{0.6em}
\textbf{Q3.} What risk appears when autoboxing meets null values?\par
\textbf{Sample answer:} Java may try to unbox null into a primitive, which throws a NullPointerException at runtime.

\end{document}
